% ============================================================
% A-19 -- Rubrik Skor Subjektif untuk Kandidat Aset Baru
% Lapisan: Skema Matematis §3 (definisi x4, x5, x6) + kontrak data v2
% Tanggal: 2026-10-02
% Status: BEKU setelah persetujuan Hafiz (lihat catatan_keputusan_20261002.md)
% ============================================================
%
% Konteks: selama ekspansi inventori aset dari 109 menuju 120 image,
% 9 dari 11 kandidat baru (gatekeeper, cilium, istio-proxy, kafka,
% velero, metrics-server, cosign, external-secrets, kube-state-metrics)
% memerlukan penilaian subjektif untuk x4 (MFA Risk), x5 (Kubernetes
% Exposure), dan x6 (Asset Criticality) karena tidak ada sumber
% telemetri otomatis untuk ketiganya pada tahap inventarisasi.
%
% Prinsip yang mendasari (lihat perdebatan "fully telemetry-driven",
% 2026-09): membuat sistem sepenuhnya berbasis telemetri untuk
% menghilangkan bias tidak menghilangkan bias -- ia memindahkan bias
% ke kode tersembunyi. Karena itu, x6 secara eksplisit adalah value
% judgment berbasis rubrik, bukan sesuatu yang diturunkan otomatis.
% Rubrik tertulis ini adalah kompromi: subjektivitas yang
% terdokumentasi dan dapat direproduksi, bukan subjektivitas yang
% tersembunyi di kepala penilai.

\subsection{Rubrik Skor Subjektif untuk Fitur x4, x5, x6 (Amendemen A-19)}
\label{subsec:a19-rubric}

Untuk aset yang belum memiliki data telemetri runtime pada saat
inventarisasi, tiga fitur dinilai menggunakan rubrik tertulis berikut.
Proses: (1) penilai menetapkan skor berdasarkan band yang paling
sesuai dengan deskripsi komponen, (2) skor dicatat beserta kutipan
band yang dipakai, (3) untuk x6, skor mentah $s \in \{0, 1, \ldots, 10\}$
dikonversi melalui $x_6 = s / 10$ sebelum masuk pipeline.

\paragraph{x6 -- Asset Criticality (skor mentah 0--10, dikonversi $x_6 = s/10$).}

\begin{itemize}
    \item \textbf{0--2 (Observational tool).} Komponen yang kegagalannya
    mengurangi visibilitas atau observabilitas tetapi tidak memengaruhi
    kerahasiaan, integritas, atau ketersediaan komponen lain secara
    langsung. Contoh: eksportir metrik read-only.

    \item \textbf{3--5 (Adjacent ke kontrol keamanan).} Komponen yang
    posisinya berdekatan dengan jalur keamanan tetapi tidak secara
    langsung menegakkan keputusan keamanan. Kompromi pada komponen ini
    dapat memfasilitasi eskalasi privilese atau manipulasi bukti secara
    tidak langsung.

    \item \textbf{6--8 (Kontrol keamanan aktif).} Komponen yang secara
    langsung menegakkan, memverifikasi, atau menyimpan state
    berkaitan-keamanan (admission control, penyimpanan secret,
    verifikasi image, penegakan kebijakan jaringan).

    \item \textbf{9--10 (Fondasional bidang-kendali klaster).}
    Komponen yang kompromi atasnya menghasilkan radius dampak
    seluruh klaster (mis. datastore klaster, API server, CNI inti).
\end{itemize}

\paragraph{x5 -- Kubernetes Exposure (langsung pada [0,1], tanpa konversi).}

\begin{itemize}
    \item \textbf{0.1--0.2 (Internal-only).} Tidak ada jalur jaringan
    langsung menuju permukaan yang terjangkau penyerang eksternal;
    komunikasi terbatas pada namespace atau ClusterIP internal.

    \item \textbf{0.3--0.4 (Dipakai luas, tidak menghadap internet).}
    Diakses oleh banyak workload lain di klaster tetapi tidak
    terekspos ke luar klaster secara langsung.

    \item \textbf{0.5 (Titik integrasi bidang-kendali/admission).}
    Disentuh oleh setiap deployment baru melalui webhook admission
    atau mekanisme serupa, menjadikannya titik ungkit luas walau
    tidak menghadap jaringan eksternal.

    \item \textbf{0.6--0.8 (Mediator lalu lintas data-plane).}
    Memediasi lalu lintas eksternal atau lintas-namespace secara
    langsung (mis. service mesh sidecar, ingress-adjacent proxy).
\end{itemize}

\paragraph{x4 -- MFA Risk (langsung pada [0,1], tanpa konversi).}

\begin{itemize}
    \item \textbf{0 (Tanpa permukaan identitas/autentikasi).}
    Komponen tidak berpartisipasi dalam alur autentikasi pengguna
    atau layanan; tidak ada kredensial yang digerbangi olehnya.

    \item \textbf{0.3 (Berdampak tidak langsung pada alur
    kredensial).} Komponen menyimpan atau meneruskan materi yang
    relevan terhadap kredensial (mis. secret) tanpa sendiri menjadi
    titik keputusan autentikasi.

    \item \textbf{0.6 (Menggerbangi keputusan autentikasi secara
    langsung).} Komponen memverifikasi atau meneruskan keputusan
    autentikasi/otorisasi sebagai bagian dari fungsinya.

    \item \textbf{1.0 (Merupakan titik penegakan MFA/autentikasi
    itu sendiri).} Komponen adalah mekanisme MFA/autentikasi utama.
\end{itemize}

\paragraph{Catatan reproduksibilitas.} Setiap skor pada tabel aset
kandidat (lih. Lampiran/\texttt{configs/assets\_candidates\_toward120.csv})
disertai kutipan band mana yang dipakai di kolom \texttt{note}, bukan
hanya angka akhir, agar audit sampel ${\ge}10\%$ (protokol mutu)
dapat memeriksa konsistensi penerapan rubrik
antar-aset, bukan hanya konsistensi angka.
